\par\noindent\textbf{(1)} 由H₀ψ=E\_nψ得 p²ψ=2m(E\_n−V)ψ，因此⟨p⁴⟩=⟨p²ψ|p²ψ⟩=4m²⟨(E\_n−V)²⟩；展开为4m²[E\_n²−2E\_n⟨V⟩+⟨V²⟩]。对Coulomb束缚态Virial定理⟨V⟩=2E\_n，故⟨p⁴⟩=4m²[−3E\_n²+⟨V²⟩]。
\par\noindent\textbf{(2)} 令精细结构常数α=e²/(4πε₀ℏc)，氢样E\_n=−mc²(Zα)²/(2n²)，且⟨r\textasciicircum{}\{-2\}⟩=Z²/[a₀²n³(l+1/2)]，a₀=ℏ/(mcα)。所以⟨V²⟩=m²c⁴(Zα)⁴/[n³(l+1/2)]，代回给 ΔE\_\{kin\}=−mc²(Zα)⁴/[2n⁴][n/(l+1/2)−3/4]。
\par\noindent\textbf{(3)} 该修正对l有显式依赖，说明纯动能相对论效应会破坏同n非相对论能级偶然简并；量级∝mc²(Zα)⁴/n³对固定l的高激发态迅速减小。推导假定Zα≪1，使速度v/c量级小，重核高Z需使用Dirac波函数。
\par\noindent\textbf{(4)} 完整电子精细结构还含自旋轨道耦合、Darwin接触项（对l=0尤重要）及相应相对论修正。它们合并后主阶只依赖总角动量j而非单独l；不可把本题ΔE\_\{kin\}误称为完整Dirac精细结构，也不能对Darwin delta项随意忽略s态贡献。
